\setcounter{section}{8}
\section{LPC17xx General Purpose Input/Output (GPIO)}

\subsection{Configuración básica}
\begin{itemize}
  \item \textbf{Power:} el capítulo indica GPIO siempre habilitado.
  \item \textbf{Pins:} el uso como GPIO depende de la selección de función realizada en la Sección 8.3.
  \item \textbf{Wake-up:} los puertos 0 y 2 pueden participar en wake-up.
  \item \textbf{Interrupts:} la habilitación por flanco se realiza con \texttt{IO0/2IntEnR} y \texttt{IO0/2IntEnF}; la IRQ se habilita además en el NVIC.
\end{itemize}

\subsection{Features}
\subsubsection{Puertos digitales}
\begin{itemize}
  \item Registros GPIO ubicados sobre bus \textbf{AHB}, con acceso rápido.
  \item Acceso por \textbf{byte}, \textbf{half-word} y \textbf{word}.
  \item Escritura del puerto completo en una instrucción.
  \item Registros separados de \texttt{SET} y \texttt{CLR}.
  \item Acceso por \textbf{GPDMA}.
  \item Soporte de \textbf{bit-banding} del Cortex-M3.
  \item Control individual de dirección por bit.
  \item Luego del reset, las E/S arrancan como entrada con \textit{pull-up}.
\end{itemize}

\paragraph{Nota técnica: bit-banding}
El bit-banding permite tratar cada bit de una región bit-bandable como una palabra alias independiente. En práctica, evita secuencias explícitas de \textit{read-modify-write} para bits individuales. En GPIO sigue siendo natural preferir \texttt{FIOSET}/\texttt{FIOCLR} para salidas puntuales, pero el soporte de bit-banding simplifica accesos bit a bit sobre registros del periférico.

\subsubsection{Puertos digitales con interrupción}
\begin{itemize}
  \item Solo \textbf{Port 0} y \textbf{Port 2} generan interrupciones GPIO por flanco.
  \item Cada pin puede generar interrupción por flanco ascendente, descendente o ambos.
  \item La detección es asíncrona, por lo que puede seguir operando sin clocks activos.
  \item Cada interrupción habilitada puede contribuir a wake-up desde \textit{Power-down}.
  \item \texttt{GPIO0} y \texttt{GPIO2} comparten posición NVIC con \texttt{EINT3}.
\end{itemize}

\subsection{Aplicaciones}
\begin{itemize}
  \item Entrada/salida digital de propósito general.
  \item Manejo de LEDs e indicadores.
  \item Control de dispositivos externos.
  \item Sensado de entradas digitales y detección de flancos.
  \item Salida de modos de bajo consumo.
\end{itemize}

\subsection{Descripción de pines}
\begin{table}[H]
  \centering
  \small
  \renewcommand{\arraystretch}{1.15}
  \setlength{\tabcolsep}{4pt}
  \caption{GPIO pin description (Table 100)}
  \begin{tabularx}{0.96\textwidth}{|>{\raggedright\arraybackslash}p{0.23\textwidth}|>{\raggedright\arraybackslash}p{0.12\textwidth}|>{\raggedright\arraybackslash}X|}
    \hline
    \textbf{Pin name} & \textbf{Type} & \textbf{Description} \\
    \hline
    \texttt{P0[30:0][1]} & Input/Output & GPIO de propósito general. Compartidos con funciones alternativas; no todos estarán disponibles en cada aplicación. \\
    \hline
    \texttt{P1[31:0][2]} & Input/Output & GPIO de propósito general. Compartidos con funciones alternativas; no todos estarán disponibles en cada aplicación. \\
    \hline
    \texttt{P2[13:0]} & Input/Output & GPIO de propósito general. \\
    \hline
    \texttt{P3[26:25]} & Input/Output & GPIO de propósito general. \\
    \hline
    \texttt{P4[29:28]} & Input/Output & GPIO de propósito general. \\
    \hline
  \end{tabularx}
\end{table}

\begin{itemize}
  \item \textbf{[1]} \texttt{P0[14:12]} no disponibles.
  \item \textbf{[2]} \texttt{P1[2]}, \texttt{P1[3]}, \texttt{P1[7:5]}, \texttt{P1[13:11]} no disponibles.
  \item Algunos pines quedan condicionados por su función alternativa. En particular, los pines de \texttt{I2C0} son \textit{open-drain} cualquiera sea la función seleccionada.
\end{itemize}

\subsection{Descripción de registros}
El bloque Fast GPIO implementa porciones de cinco puertos de 32 bits. Los registros \texttt{FIOx*} residen en \texttt{0x2009 C000} a \texttt{0x2009 C09C} sobre AHB. Los registros de interrupción GPIO residen en \texttt{0x4002 8080} a \texttt{0x4002 80B4}.

\renewcommand{\arraystretch}{1.12}
\setlength{\LTleft}{0pt}
\setlength{\LTright}{0pt}

\begin{longtable}{|>{\raggedright\arraybackslash}p{0.14\textwidth}|>{\raggedright\arraybackslash}p{0.34\textwidth}|>{\centering\arraybackslash}p{0.07\textwidth}|>{\centering\arraybackslash}p{0.08\textwidth}|>{\raggedright\arraybackslash}p{0.23\textwidth}|}
\caption{GPIO register map (local bus accessible registers - enhanced GPIO features) (Table 101)} \\
\hline
\textbf{Name} & \textbf{Description} & \textbf{Access} & \textbf{Reset} & \textbf{PORTn Register Name \& Address} \\
\hline
\endfirsthead
\hline
\textbf{Name} & \textbf{Description} & \textbf{Access} & \textbf{Reset} & \textbf{PORTn Register Name \& Address} \\
\hline
\endhead
\texttt{FIODIR} & Control de dirección por bit del puerto. & R/W & 0 & \texttt{FIO0DIR - 0x2009 C000}, \texttt{FIO1DIR - 0x2009 C020}, \texttt{FIO2DIR - 0x2009 C040}, \texttt{FIO3DIR - 0x2009 C060}, \texttt{FIO4DIR - 0x2009 C080}. \\
\hline
\texttt{FIOMASK} & Máscara del puerto. Las lecturas y escrituras por \texttt{FIOPIN}, \texttt{FIOSET} y \texttt{FIOCLR} solo afectan bits habilitados por ceros en esta máscara. & R/W & 0 & \texttt{FIO0MASK - 0x2009 C010}, \texttt{FIO1MASK - 0x2009 C030}, \texttt{FIO2MASK - 0x2009 C050}, \texttt{FIO3MASK - 0x2009 C070}, \texttt{FIO4MASK - 0x2009 C090}. \\
\hline
\texttt{FIOPIN} & Valor del pin usando \texttt{FIOMASK}. Puede leerse independientemente de la dirección o de la función digital alternativa; al escribir, aplica valores sobre bits no enmascarados. & R/W & 0 & \texttt{FIO0PIN - 0x2009 C014}, \texttt{FIO1PIN - 0x2009 C034}, \texttt{FIO2PIN - 0x2009 C054}, \texttt{FIO3PIN - 0x2009 C074}, \texttt{FIO4PIN - 0x2009 C094}. \\
\hline
\texttt{FIOSET} & Fuerza niveles altos en pines de salida no enmascarados. Escribir \texttt{1} seta; escribir \texttt{0} no altera. & R/W & 0 & \texttt{FIO0SET - 0x2009 C018}, \texttt{FIO1SET - 0x2009 C038}, \texttt{FIO2SET - 0x2009 C058}, \texttt{FIO3SET - 0x2009 C078}, \texttt{FIO4SET - 0x2009 C098}. \\
\hline
\texttt{FIOCLR} & Fuerza niveles bajos en pines de salida no enmascarados. Escribir \texttt{1} limpia; escribir \texttt{0} no altera. & WO & 0 & \texttt{FIO0CLR - 0x2009 C01C}, \texttt{FIO1CLR - 0x2009 C03C}, \texttt{FIO2CLR - 0x2009 C05C}, \texttt{FIO3CLR - 0x2009 C07C}, \texttt{FIO4CLR - 0x2009 C09C}. \\
\hline
\end{longtable}

\begin{table}[H]
  \centering
  \small
  \setlength{\tabcolsep}{4pt}
  \caption{GPIO interrupt register map (Table 102)}
  \begin{tabularx}{0.96\textwidth}{|>{\raggedright\arraybackslash}p{0.24\textwidth}|>{\raggedright\arraybackslash}X|>{\centering\arraybackslash}p{0.11\textwidth}|>{\centering\arraybackslash}p{0.11\textwidth}|}
    \hline
    \textbf{Name} & \textbf{Description} & \textbf{Access} & \textbf{Reset} \\
    \hline
    \texttt{IO0IntEnR / IO2IntEnR} & GPIO Interrupt Enable for Rising edge. & R/W & 0 \\
    \hline
    \texttt{IO0IntEnF / IO2IntEnF} & GPIO Interrupt Enable for Falling edge. & R/W & 0 \\
    \hline
    \texttt{IO0IntStatR / IO2IntStatR} & GPIO Interrupt Status for Rising edge. & RO & 0 \\
    \hline
    \texttt{IO0IntStatF / IO2IntStatF} & GPIO Interrupt Status for Falling edge. & RO & 0 \\
    \hline
    \texttt{IO0IntClr / IO2IntClr} & GPIO Interrupt Clear. & WO & 0 \\
    \hline
    \texttt{IOIntStatus} & GPIO overall Interrupt Status. & RO & 0 \\
    \hline
  \end{tabularx}
\end{table}

\subsubsection{GPIO port Direction register \texttt{FIOxDIR}}
\begin{itemize}
  \item Controla la dirección de los pines configurados como GPIO.
  \item \texttt{0}: entrada.
  \item \texttt{1}: salida.
  \item \texttt{P0.29} y \texttt{P0.30} comparten pad con \texttt{USB\_D+} y \texttt{USB\_D-}.
\end{itemize}

\begin{table}[H]
  \centering
  \small
  \caption{Fast GPIO port Direction register \texttt{FIO0DIR} to \texttt{FIO4DIR} (Table 103)}
  \begin{tabularx}{0.96\textwidth}{|>{\centering\arraybackslash}p{0.10\textwidth}|>{\raggedright\arraybackslash}p{0.19\textwidth}|>{\raggedright\arraybackslash}X|>{\centering\arraybackslash}p{0.10\textwidth}|}
    \hline
    \textbf{Bit} & \textbf{Symbol} & \textbf{Value / Description} & \textbf{Reset} \\
    \hline
    \texttt{31:0} & \texttt{FIOxDIR} & Bit 0 controla \texttt{Px.0}; bit 31 controla \texttt{Px.31}. \texttt{0}: pin configurado como entrada. \texttt{1}: pin configurado como salida. & \texttt{0x0} \\
    \hline
  \end{tabularx}
\end{table}

\begin{longtable}{|>{\raggedright\arraybackslash}p{0.15\textwidth}|>{\raggedright\arraybackslash}p{0.35\textwidth}|>{\centering\arraybackslash}p{0.10\textwidth}|>{\centering\arraybackslash}p{0.08\textwidth}|>{\raggedright\arraybackslash}p{0.20\textwidth}|}
\caption{Fast GPIO port Direction control byte and half-word accessible register description (Table 104)} \\
\hline
\textbf{Generic register} & \textbf{Description} & \textbf{Length / Access} & \textbf{Reset} & \textbf{PORTn Register Address} \\
\hline
\endfirsthead
\hline
\textbf{Generic register} & \textbf{Description} & \textbf{Length / Access} & \textbf{Reset} & \textbf{PORTn Register Address} \\
\hline
\endhead
\texttt{FIOxDIR0} & Control de dirección para \texttt{Px.0} a \texttt{Px.7}. & 8 / R/W & \texttt{0x00} & \texttt{FIO0DIR0 - 0x2009 C000}, \texttt{FIO1DIR0 - 0x2009 C020}, \texttt{FIO2DIR0 - 0x2009 C040}, \texttt{FIO3DIR0 - 0x2009 C060}, \texttt{FIO4DIR0 - 0x2009 C080}. \\
\hline
\texttt{FIOxDIR1} & Control de dirección para \texttt{Px.8} a \texttt{Px.15}. & 8 / R/W & \texttt{0x00} & \texttt{FIO0DIR1 - 0x2009 C001}, \texttt{FIO1DIR1 - 0x2009 C021}, \texttt{FIO2DIR1 - 0x2009 C041}, \texttt{FIO3DIR1 - 0x2009 C061}, \texttt{FIO4DIR1 - 0x2009 C081}. \\
\hline
\texttt{FIOxDIR2} & Control de dirección para \texttt{Px.16} a \texttt{Px.23}. & 8 / R/W & \texttt{0x00} & \texttt{FIO0DIR2 - 0x2009 C002}, \texttt{FIO1DIR2 - 0x2009 C022}, \texttt{FIO2DIR2 - 0x2009 C042}, \texttt{FIO3DIR2 - 0x2009 C062}, \texttt{FIO4DIR2 - 0x2009 C082}. \\
\hline
\texttt{FIOxDIR3} & Control de dirección para \texttt{Px.24} a \texttt{Px.31}. & 8 / R/W & \texttt{0x00} & \texttt{FIO0DIR3 - 0x2009 C003}, \texttt{FIO1DIR3 - 0x2009 C023}, \texttt{FIO2DIR3 - 0x2009 C043}, \texttt{FIO3DIR3 - 0x2009 C063}, \texttt{FIO4DIR3 - 0x2009 C083}. \\
\hline
\texttt{FIOxDIRL} & Control de dirección para la mitad baja del puerto, \texttt{Px.0} a \texttt{Px.15}. & 16 / R/W & \texttt{0x0000} & \texttt{FIO0DIRL - 0x2009 C000}, \texttt{FIO1DIRL - 0x2009 C020}, \texttt{FIO2DIRL - 0x2009 C040}, \texttt{FIO3DIRL - 0x2009 C060}, \texttt{FIO4DIRL - 0x2009 C080}. \\
\hline
\texttt{FIOxDIRU} & Control de dirección para la mitad alta del puerto, \texttt{Px.16} a \texttt{Px.31}. & 16 / R/W & \texttt{0x0000} & \texttt{FIO0DIRU - 0x2009 C002}, \texttt{FIO1DIRU - 0x2009 C022}, \texttt{FIO2DIRU - 0x2009 C042}, \texttt{FIO3DIRU - 0x2009 C062}, \texttt{FIO4DIRU - 0x2009 C082}. \\
\hline
\end{longtable}

\subsubsection{GPIO port output Set register \texttt{FIOxSET}}
\begin{itemize}
  \item Escribir \texttt{1} fuerza nivel alto sobre el pin de salida correspondiente.
  \item Escribir \texttt{0} no altera el estado actual.
  \item La operación queda condicionada por la máscara \texttt{FIOxMASK}.
\end{itemize}

\begin{table}[H]
  \centering
  \small
  \caption{Fast GPIO port output Set register (Table 105)}
  \begin{tabularx}{0.96\textwidth}{|>{\centering\arraybackslash}p{0.10\textwidth}|>{\raggedright\arraybackslash}p{0.19\textwidth}|>{\raggedright\arraybackslash}X|>{\centering\arraybackslash}p{0.10\textwidth}|}
    \hline
    \textbf{Bit} & \textbf{Symbol} & \textbf{Value / Description} & \textbf{Reset} \\
    \hline
    \texttt{31:0} & \texttt{FIOxSET} & Bit 0 controla \texttt{Px.0}; bit 31 controla \texttt{Px.31}. \texttt{0}: salida sin cambios. \texttt{1}: salida forzada a HIGH. & \texttt{0x0} \\
    \hline
  \end{tabularx}
\end{table}

\begin{longtable}{|>{\raggedright\arraybackslash}p{0.15\textwidth}|>{\raggedright\arraybackslash}p{0.35\textwidth}|>{\centering\arraybackslash}p{0.10\textwidth}|>{\centering\arraybackslash}p{0.08\textwidth}|>{\raggedright\arraybackslash}p{0.20\textwidth}|}
\caption{Fast GPIO port output Set byte and half-word accessible register description (Table 106)} \\
\hline
\textbf{Generic register} & \textbf{Description} & \textbf{Length / Access} & \textbf{Reset} & \textbf{PORTn Register Address} \\
\hline
\endfirsthead
\hline
\textbf{Generic register} & \textbf{Description} & \textbf{Length / Access} & \textbf{Reset} & \textbf{PORTn Register Address} \\
\hline
\endhead
\texttt{FIOxSET0} & Set de salida para \texttt{Px.0} a \texttt{Px.7}. & 8 / R/W & \texttt{0x00} & \texttt{FIO0SET0 - 0x2009 C018}, \texttt{FIO1SET0 - 0x2009 C038}, \texttt{FIO2SET0 - 0x2009 C058}, \texttt{FIO3SET0 - 0x2009 C078}, \texttt{FIO4SET0 - 0x2009 C098}. \\
\hline
\texttt{FIOxSET1} & Set de salida para \texttt{Px.8} a \texttt{Px.15}. & 8 / R/W & \texttt{0x00} & \texttt{FIO0SET1 - 0x2009 C019}, \texttt{FIO1SET1 - 0x2009 C039}, \texttt{FIO2SET1 - 0x2009 C059}, \texttt{FIO3SET1 - 0x2009 C079}, \texttt{FIO4SET1 - 0x2009 C099}. \\
\hline
\texttt{FIOxSET2} & Set de salida para \texttt{Px.16} a \texttt{Px.23}. & 8 / R/W & \texttt{0x00} & \texttt{FIO0SET2 - 0x2009 C01A}, \texttt{FIO1SET2 - 0x2009 C03A}, \texttt{FIO2SET2 - 0x2009 C05A}, \texttt{FIO3SET2 - 0x2009 C07A}, \texttt{FIO4SET2 - 0x2009 C09A}. \\
\hline
\texttt{FIOxSET3} & Set de salida para \texttt{Px.24} a \texttt{Px.31}. & 8 / R/W & \texttt{0x00} & \texttt{FIO0SET3 - 0x2009 C01B}, \texttt{FIO1SET3 - 0x2009 C03B}, \texttt{FIO2SET3 - 0x2009 C05B}, \texttt{FIO3SET3 - 0x2009 C07B}, \texttt{FIO4SET3 - 0x2009 C09B}. \\
\hline
\texttt{FIOxSETL} & Set de salida para \texttt{Px.0} a \texttt{Px.15}. & 16 / R/W & \texttt{0x0000} & \texttt{FIO0SETL - 0x2009 C018}, \texttt{FIO1SETL - 0x2009 C038}, \texttt{FIO2SETL - 0x2009 C058}, \texttt{FIO3SETL - 0x2009 C078}, \texttt{FIO4SETL - 0x2009 C098}. \\
\hline
\texttt{FIOxSETU} & Set de salida para \texttt{Px.16} a \texttt{Px.31}. & 16 / R/W & \texttt{0x0000} & \texttt{FIO0SETU - 0x2009 C01A}, \texttt{FIO1SETU - 0x2009 C03A}, \texttt{FIO2SETU - 0x2009 C05A}, \texttt{FIO3SETU - 0x2009 C07A}, \texttt{FIO4SETU - 0x2009 C09A}. \\
\hline
\end{longtable}

\subsubsection{GPIO port output Clear register \texttt{FIOxCLR}}
\begin{itemize}
  \item Escribir \texttt{1} fuerza nivel bajo sobre el pin de salida correspondiente.
  \item Escribir \texttt{0} no altera el estado actual.
\end{itemize}

\begin{table}[H]
  \centering
  \small
  \caption{Fast GPIO port output Clear register (Table 107)}
  \begin{tabularx}{0.96\textwidth}{|>{\centering\arraybackslash}p{0.10\textwidth}|>{\raggedright\arraybackslash}p{0.19\textwidth}|>{\raggedright\arraybackslash}X|>{\centering\arraybackslash}p{0.10\textwidth}|}
    \hline
    \textbf{Bit} & \textbf{Symbol} & \textbf{Value / Description} & \textbf{Reset} \\
    \hline
    \texttt{31:0} & \texttt{FIOxCLR} & Bit 0 controla \texttt{Px.0}; bit 31 controla \texttt{Px.31}. \texttt{0}: salida sin cambios. \texttt{1}: salida forzada a LOW. & \texttt{0x0} \\
    \hline
  \end{tabularx}
\end{table}

\begin{longtable}{|>{\raggedright\arraybackslash}p{0.15\textwidth}|>{\raggedright\arraybackslash}p{0.35\textwidth}|>{\centering\arraybackslash}p{0.10\textwidth}|>{\centering\arraybackslash}p{0.08\textwidth}|>{\raggedright\arraybackslash}p{0.20\textwidth}|}
\caption{Fast GPIO port output Clear byte and half-word accessible register description (Table 108)} \\
\hline
\textbf{Generic register} & \textbf{Description} & \textbf{Length / Access} & \textbf{Reset} & \textbf{PORTn Register Address} \\
\hline
\endfirsthead
\hline
\textbf{Generic register} & \textbf{Description} & \textbf{Length / Access} & \textbf{Reset} & \textbf{PORTn Register Address} \\
\hline
\endhead
\texttt{FIOxCLR0} & Clear de salida para \texttt{Px.0} a \texttt{Px.7}. & 8 / R/W & \texttt{0x00} & \texttt{FIO0CLR0 - 0x2009 C01C}, \texttt{FIO1CLR0 - 0x2009 C03C}, \texttt{FIO2CLR0 - 0x2009 C05C}, \texttt{FIO3CLR0 - 0x2009 C07C}, \texttt{FIO4CLR0 - 0x2009 C09C}. \\
\hline
\texttt{FIOxCLR1} & Clear de salida para \texttt{Px.8} a \texttt{Px.15}. & 8 / R/W & \texttt{0x00} & \texttt{FIO0CLR1 - 0x2009 C01D}, \texttt{FIO1CLR1 - 0x2009 C03D}, \texttt{FIO2CLR1 - 0x2009 C05D}, \texttt{FIO3CLR1 - 0x2009 C07D}, \texttt{FIO4CLR1 - 0x2009 C09D}. \\
\hline
\texttt{FIOxCLR2} & Clear de salida para \texttt{Px.16} a \texttt{Px.23}. & 8 / R/W & \texttt{0x00} & \texttt{FIO0CLR2 - 0x2009 C01E}, \texttt{FIO1CLR2 - 0x2009 C03E}, \texttt{FIO2CLR2 - 0x2009 C05E}, \texttt{FIO3CLR2 - 0x2009 C07E}, \texttt{FIO4CLR2 - 0x2009 C09E}. \\
\hline
\texttt{FIOxCLR3} & Clear de salida para \texttt{Px.24} a \texttt{Px.31}. & 8 / R/W & \texttt{0x00} & \texttt{FIO0CLR3 - 0x2009 C01F}, \texttt{FIO1CLR3 - 0x2009 C03F}, \texttt{FIO2CLR3 - 0x2009 C05F}, \texttt{FIO3CLR3 - 0x2009 C07F}, \texttt{FIO4CLR3 - 0x2009 C09F}. \\
\hline
\texttt{FIOxCLRL} & Clear de salida para \texttt{Px.0} a \texttt{Px.15}. & 16 / R/W & \texttt{0x0000} & \texttt{FIO0CLRL - 0x2009 C01C}, \texttt{FIO1CLRL - 0x2009 C03C}, \texttt{FIO2CLRL - 0x2009 C05C}, \texttt{FIO3CLRL - 0x2009 C07C}, \texttt{FIO4CLRL - 0x2009 C09C}. \\
\hline
\texttt{FIOxCLRU} & Clear de salida para \texttt{Px.16} a \texttt{Px.31}. & 16 / R/W & \texttt{0x0000} & \texttt{FIO0CLRU - 0x2009 C01E}, \texttt{FIO1CLRU - 0x2009 C03E}, \texttt{FIO2CLRU - 0x2009 C05E}, \texttt{FIO3CLRU - 0x2009 C07E}, \texttt{FIO4CLRU - 0x2009 C09E}. \\
\hline
\end{longtable}

\subsubsection{GPIO port Pin value register \texttt{FIOxPIN}}
\begin{itemize}
  \item Permite leer el estado digital actual del pin, tanto si el pin está como entrada como si está como salida.
  \item También puede escribir valores sobre bits no enmascarados.
  \item Si se lee \texttt{FIOPIN}, los bits enmascarados con \texttt{1} en \texttt{FIOMASK} retornan \texttt{0} aun cuando el pin físico esté en otro estado.
\end{itemize}

\begin{table}[H]
  \centering
  \small
  \caption{Fast GPIO port Pin value register (Table 109)}
  \begin{tabularx}{0.96\textwidth}{|>{\centering\arraybackslash}p{0.10\textwidth}|>{\raggedright\arraybackslash}p{0.19\textwidth}|>{\raggedright\arraybackslash}X|>{\centering\arraybackslash}p{0.10\textwidth}|}
    \hline
    \textbf{Bit} & \textbf{Symbol} & \textbf{Value / Description} & \textbf{Reset} \\
    \hline
    \texttt{31:0} & \texttt{FIOxPIN} & Valor del pin \texttt{Px.n}. La lectura refleja el estado digital actual. La escritura aplica el valor escrito sobre bits no enmascarados del puerto. & \texttt{0x0} \\
    \hline
  \end{tabularx}
\end{table}

\begin{longtable}{|>{\raggedright\arraybackslash}p{0.15\textwidth}|>{\raggedright\arraybackslash}p{0.35\textwidth}|>{\centering\arraybackslash}p{0.10\textwidth}|>{\centering\arraybackslash}p{0.08\textwidth}|>{\raggedright\arraybackslash}p{0.20\textwidth}|}
\caption{Fast GPIO port Pin value byte and half-word accessible register description (Table 110)} \\
\hline
\textbf{Generic register} & \textbf{Description} & \textbf{Length / Access} & \textbf{Reset} & \textbf{PORTn Register Address} \\
\hline
\endfirsthead
\hline
\textbf{Generic register} & \textbf{Description} & \textbf{Length / Access} & \textbf{Reset} & \textbf{PORTn Register Address} \\
\hline
\endhead
\texttt{FIOxPIN0} & Valor del pin para \texttt{Px.0} a \texttt{Px.7}. & 8 / R/W & \texttt{0x00} & \texttt{FIO0PIN0 - 0x2009 C014}, \texttt{FIO1PIN0 - 0x2009 C034}, \texttt{FIO2PIN0 - 0x2009 C054}, \texttt{FIO3PIN0 - 0x2009 C074}, \texttt{FIO4PIN0 - 0x2009 C094}. \\
\hline
\texttt{FIOxPIN1} & Valor del pin para \texttt{Px.8} a \texttt{Px.15}. & 8 / R/W & \texttt{0x00} & \texttt{FIO0PIN1 - 0x2009 C015}, \texttt{FIO1PIN1 - 0x2009 C035}, \texttt{FIO2PIN1 - 0x2009 C055}, \texttt{FIO3PIN1 - 0x2009 C075}, \texttt{FIO4PIN1 - 0x2009 C095}. \\
\hline
\texttt{FIOxPIN2} & Valor del pin para \texttt{Px.16} a \texttt{Px.23}. & 8 / R/W & \texttt{0x00} & \texttt{FIO0PIN2 - 0x2009 C016}, \texttt{FIO1PIN2 - 0x2009 C036}, \texttt{FIO2PIN2 - 0x2009 C056}, \texttt{FIO3PIN2 - 0x2009 C076}, \texttt{FIO4PIN2 - 0x2009 C096}. \\
\hline
\texttt{FIOxPIN3} & Valor del pin para \texttt{Px.24} a \texttt{Px.31}. & 8 / R/W & \texttt{0x00} & \texttt{FIO0PIN3 - 0x2009 C017}, \texttt{FIO1PIN3 - 0x2009 C037}, \texttt{FIO2PIN3 - 0x2009 C057}, \texttt{FIO3PIN3 - 0x2009 C077}, \texttt{FIO4PIN3 - 0x2009 C097}. \\
\hline
\texttt{FIOxPINL} & Valor del pin para \texttt{Px.0} a \texttt{Px.15}. & 16 / R/W & \texttt{0x0000} & \texttt{FIO0PINL - 0x2009 C014}, \texttt{FIO1PINL - 0x2009 C034}, \texttt{FIO2PINL - 0x2009 C054}, \texttt{FIO3PINL - 0x2009 C074}, \texttt{FIO4PINL - 0x2009 C094}. \\
\hline
\texttt{FIOxPINU} & Valor del pin para \texttt{Px.16} a \texttt{Px.31}. & 16 / R/W & \texttt{0x0000} & \texttt{FIO0PINU - 0x2009 C016}, \texttt{FIO1PINU - 0x2009 C036}, \texttt{FIO2PINU - 0x2009 C056}, \texttt{FIO3PINU - 0x2009 C076}, \texttt{FIO4PINU - 0x2009 C096}. \\
\hline
\end{longtable}

\subsubsection{Fast GPIO port Mask register \texttt{FIOxMASK}}
\begin{itemize}
  \item Selecciona qué bits participan en lecturas y escrituras del puerto.
  \item \texttt{0}: bit habilitado.
  \item \texttt{1}: bit enmascarado.
\end{itemize}

\begin{table}[H]
  \centering
  \small
  \caption{Fast GPIO port Mask register (Table 111)}
  \begin{tabularx}{0.96\textwidth}{|>{\centering\arraybackslash}p{0.10\textwidth}|>{\raggedright\arraybackslash}p{0.19\textwidth}|>{\raggedright\arraybackslash}X|>{\centering\arraybackslash}p{0.10\textwidth}|}
    \hline
    \textbf{Bit} & \textbf{Symbol} & \textbf{Value / Description} & \textbf{Reset} \\
    \hline
    \texttt{31:0} & \texttt{FIOxMASK} & \texttt{0}: el bit físico participa de lecturas/escrituras. \texttt{1}: el bit queda enmascarado; lecturas retornan \texttt{0} y escrituras no alteran el pin. & \texttt{0x0} \\
    \hline
  \end{tabularx}
\end{table}

\begin{longtable}{|>{\raggedright\arraybackslash}p{0.15\textwidth}|>{\raggedright\arraybackslash}p{0.35\textwidth}|>{\centering\arraybackslash}p{0.10\textwidth}|>{\centering\arraybackslash}p{0.08\textwidth}|>{\raggedright\arraybackslash}p{0.20\textwidth}|}
\caption{Fast GPIO port Mask byte and half-word accessible register description (Table 112)} \\
\hline
\textbf{Generic register} & \textbf{Description} & \textbf{Length / Access} & \textbf{Reset} & \textbf{PORTn Register Address} \\
\hline
\endfirsthead
\hline
\textbf{Generic register} & \textbf{Description} & \textbf{Length / Access} & \textbf{Reset} & \textbf{PORTn Register Address} \\
\hline
\endhead
\texttt{FIOxMASK0} & Máscara para \texttt{Px.0} a \texttt{Px.7}. & 8 / R/W & \texttt{0x00} & \texttt{FIO0MASK0 - 0x2009 C010}, \texttt{FIO1MASK0 - 0x2009 C030}, \texttt{FIO2MASK0 - 0x2009 C050}, \texttt{FIO3MASK0 - 0x2009 C070}, \texttt{FIO4MASK0 - 0x2009 C090}. \\
\hline
\texttt{FIOxMASK1} & Máscara para \texttt{Px.8} a \texttt{Px.15}. & 8 / R/W & \texttt{0x00} & \texttt{FIO0MASK1 - 0x2009 C011}, \texttt{FIO1MASK1 - 0x2009 C031}, \texttt{FIO2MASK1 - 0x2009 C051}, \texttt{FIO3MASK1 - 0x2009 C071}, \texttt{FIO4MASK1 - 0x2009 C091}. \\
\hline
\texttt{FIOxMASK2} & Máscara para \texttt{Px.16} a \texttt{Px.23}. & 8 / R/W & \texttt{0x00} & \texttt{FIO0MASK2 - 0x2009 C012}, \texttt{FIO1MASK2 - 0x2009 C032}, \texttt{FIO2MASK2 - 0x2009 C052}, \texttt{FIO3MASK2 - 0x2009 C072}, \texttt{FIO4MASK2 - 0x2009 C092}. \\
\hline
\texttt{FIOxMASK3} & Máscara para \texttt{Px.24} a \texttt{Px.31}. & 8 / R/W & \texttt{0x00} & \texttt{FIO0MASK3 - 0x2009 C013}, \texttt{FIO1MASK3 - 0x2009 C033}, \texttt{FIO2MASK3 - 0x2009 C053}, \texttt{FIO3MASK3 - 0x2009 C073}, \texttt{FIO4MASK3 - 0x2009 C093}. \\
\hline
\texttt{FIOxMASKL} & Máscara para \texttt{Px.0} a \texttt{Px.15}. & 16 / R/W & \texttt{0x0000} & \texttt{FIO0MASKL - 0x2009 C010}, \texttt{FIO1MASKL - 0x2009 C030}, \texttt{FIO2MASKL - 0x2009 C050}, \texttt{FIO3MASKL - 0x2009 C070}, \texttt{FIO4MASKL - 0x2009 C090}. \\
\hline
\texttt{FIOxMASKU} & Máscara para \texttt{Px.16} a \texttt{Px.31}. & 16 / R/W & \texttt{0x0000} & \texttt{FIO0MASKU - 0x2009 C012}, \texttt{FIO1MASKU - 0x2009 C032}, \texttt{FIO2MASKU - 0x2009 C052}, \texttt{FIO3MASKU - 0x2009 C072}, \texttt{FIO4MASKU - 0x2009 C092}. \\
\hline
\end{longtable}

\subsubsection{GPIO interrupt registers}

\begin{table}[H]
  \centering
  \small
  \caption{GPIO overall Interrupt Status register \texttt{IOIntStatus} (Table 113)}
  \begin{tabularx}{0.96\textwidth}{|>{\centering\arraybackslash}p{0.12\textwidth}|>{\raggedright\arraybackslash}p{0.18\textwidth}|>{\raggedright\arraybackslash}X|>{\centering\arraybackslash}p{0.10\textwidth}|}
    \hline
    \textbf{Bit} & \textbf{Symbol} & \textbf{Value / Description} & \textbf{Reset} \\
    \hline
    \texttt{0} & \texttt{P0Int} & \texttt{0}: no hay interrupciones pendientes en Port 0. \texttt{1}: existe al menos una interrupción pendiente en Port 0. & \texttt{0} \\
    \hline
    \texttt{1} & \texttt{-} & Reservado. & NA \\
    \hline
    \texttt{2} & \texttt{P2Int} & \texttt{0}: no hay interrupciones pendientes en Port 2. \texttt{1}: existe al menos una interrupción pendiente en Port 2. & \texttt{0} \\
    \hline
    \texttt{31:3} & \texttt{-} & Reservado. & NA \\
    \hline
  \end{tabularx}
\end{table}

\begin{table}[H]
  \centering
  \small
  \caption{GPIO Interrupt Enable for port 0 Rising Edge \texttt{IO0IntEnR} (Table 114)}
  \begin{tabularx}{0.96\textwidth}{|>{\centering\arraybackslash}p{0.14\textwidth}|>{\raggedright\arraybackslash}p{0.18\textwidth}|>{\raggedright\arraybackslash}X|>{\centering\arraybackslash}p{0.10\textwidth}|}
    \hline
    \textbf{Bit} & \textbf{Symbol} & \textbf{Value / Description} & \textbf{Reset} \\
    \hline
    \texttt{0--11, 15--18} & \texttt{P0.nER} & Habilita interrupción por flanco ascendente para \texttt{P0.n}. \texttt{0}: deshabilitada. \texttt{1}: habilitada. & \texttt{0} \\
    \hline
    \texttt{19--30[1]} & \texttt{P0.nER} & Igual que la fila anterior para los bits altos disponibles en Port 0. & \texttt{0} \\
    \hline
    \texttt{14:12} & \texttt{-} & Reservado. & NA \\
    \hline
    \texttt{31} & \texttt{-} & Reservado. & NA \\
    \hline
  \end{tabularx}
\end{table}

\begin{table}[H]
  \centering
  \small
  \caption{GPIO Interrupt Enable for port 2 Rising Edge \texttt{IO2IntEnR} (Table 115)}
  \begin{tabularx}{0.96\textwidth}{|>{\centering\arraybackslash}p{0.14\textwidth}|>{\raggedright\arraybackslash}p{0.18\textwidth}|>{\raggedright\arraybackslash}X|>{\centering\arraybackslash}p{0.10\textwidth}|}
    \hline
    \textbf{Bit} & \textbf{Symbol} & \textbf{Value / Description} & \textbf{Reset} \\
    \hline
    \texttt{0--11} & \texttt{P2.nER} & Habilita interrupción por flanco ascendente para \texttt{P2.n}. \texttt{0}: deshabilitada. \texttt{1}: habilitada. & \texttt{0} \\
    \hline
    \texttt{12--13[1]} & \texttt{P2.nER} & Igual que la fila anterior para \texttt{P2.12} y \texttt{P2.13}. & \texttt{0} \\
    \hline
    \texttt{31:14} & \texttt{-} & Reservado. & NA \\
    \hline
  \end{tabularx}
\end{table}

\begin{table}[H]
  \centering
  \small
  \caption{GPIO Interrupt Enable for port 0 Falling Edge \texttt{IO0IntEnF} (Table 116)}
  \begin{tabularx}{0.96\textwidth}{|>{\centering\arraybackslash}p{0.14\textwidth}|>{\raggedright\arraybackslash}p{0.18\textwidth}|>{\raggedright\arraybackslash}X|>{\centering\arraybackslash}p{0.10\textwidth}|}
    \hline
    \textbf{Bit} & \textbf{Symbol} & \textbf{Value / Description} & \textbf{Reset} \\
    \hline
    \texttt{0--11, 15--18} & \texttt{P0.nEF} & Habilita interrupción por flanco descendente para \texttt{P0.n}. \texttt{0}: deshabilitada. \texttt{1}: habilitada. & \texttt{0} \\
    \hline
    \texttt{19--30[1]} & \texttt{P0.nEF} & Igual que la fila anterior para los bits altos disponibles en Port 0. & \texttt{0} \\
    \hline
    \texttt{14:12} & \texttt{-} & Reservado. & NA \\
    \hline
    \texttt{31} & \texttt{-} & Reservado. & NA \\
    \hline
  \end{tabularx}
\end{table}

\begin{table}[H]
  \centering
  \small
  \caption{GPIO Interrupt Enable for port 2 Falling Edge \texttt{IO2IntEnF} (Table 117)}
  \begin{tabularx}{0.96\textwidth}{|>{\centering\arraybackslash}p{0.14\textwidth}|>{\raggedright\arraybackslash}p{0.18\textwidth}|>{\raggedright\arraybackslash}X|>{\centering\arraybackslash}p{0.10\textwidth}|}
    \hline
    \textbf{Bit} & \textbf{Symbol} & \textbf{Value / Description} & \textbf{Reset} \\
    \hline
    \texttt{0--11} & \texttt{P2.nEF} & Habilita interrupción por flanco descendente para \texttt{P2.n}. \texttt{0}: deshabilitada. \texttt{1}: habilitada. & \texttt{0} \\
    \hline
    \texttt{12--13[1]} & \texttt{P2.nEF} & Igual que la fila anterior para \texttt{P2.12} y \texttt{P2.13}. & \texttt{0} \\
    \hline
    \texttt{31:14} & \texttt{-} & Reservado. & NA \\
    \hline
  \end{tabularx}
\end{table}

\begin{table}[H]
  \centering
  \small
  \caption{GPIO Interrupt Status for port 0 Rising Edge Interrupt \texttt{IO0IntStatR} (Table 118)}
  \begin{tabularx}{0.96\textwidth}{|>{\centering\arraybackslash}p{0.14\textwidth}|>{\raggedright\arraybackslash}p{0.18\textwidth}|>{\raggedright\arraybackslash}X|>{\centering\arraybackslash}p{0.10\textwidth}|}
    \hline
    \textbf{Bit} & \textbf{Symbol} & \textbf{Value / Description} & \textbf{Reset} \\
    \hline
    \texttt{0--11, 15--18} & \texttt{P0.nREI} & Estado de interrupción por flanco ascendente en \texttt{P0.n}. \texttt{0}: no detectado. \texttt{1}: interrupción generada. & \texttt{0} \\
    \hline
    \texttt{19--30[1]} & \texttt{P0.nREI} & Igual que la fila anterior para los bits altos disponibles en Port 0. & \texttt{0} \\
    \hline
    \texttt{14:12} & \texttt{-} & Reservado. & NA \\
    \hline
    \texttt{31} & \texttt{-} & Reservado. & NA \\
    \hline
  \end{tabularx}
\end{table}

\begin{table}[H]
  \centering
  \small
  \caption{GPIO Interrupt Status for port 2 Rising Edge Interrupt \texttt{IO2IntStatR} (Table 119)}
  \begin{tabularx}{0.96\textwidth}{|>{\centering\arraybackslash}p{0.14\textwidth}|>{\raggedright\arraybackslash}p{0.18\textwidth}|>{\raggedright\arraybackslash}X|>{\centering\arraybackslash}p{0.10\textwidth}|}
    \hline
    \textbf{Bit} & \textbf{Symbol} & \textbf{Value / Description} & \textbf{Reset} \\
    \hline
    \texttt{0--11} & \texttt{P2.nREI} & Estado de interrupción por flanco ascendente en \texttt{P2.n}. \texttt{0}: no detectado. \texttt{1}: interrupción generada. & \texttt{0} \\
    \hline
    \texttt{12--13[1]} & \texttt{P2.nREI} & Igual que la fila anterior para \texttt{P2.12} y \texttt{P2.13}. & \texttt{0} \\
    \hline
    \texttt{31:14} & \texttt{-} & Reservado. & NA \\
    \hline
  \end{tabularx}
\end{table}

\begin{table}[H]
  \centering
  \small
  \caption{GPIO Interrupt Status for port 0 Falling Edge Interrupt \texttt{IO0IntStatF} (Table 120)}
  \begin{tabularx}{0.96\textwidth}{|>{\centering\arraybackslash}p{0.14\textwidth}|>{\raggedright\arraybackslash}p{0.18\textwidth}|>{\raggedright\arraybackslash}X|>{\centering\arraybackslash}p{0.10\textwidth}|}
    \hline
    \textbf{Bit} & \textbf{Symbol} & \textbf{Value / Description} & \textbf{Reset} \\
    \hline
    \texttt{0--11, 15--18} & \texttt{P0.nFEI} & Estado de interrupción por flanco descendente en \texttt{P0.n}. \texttt{0}: no detectado. \texttt{1}: interrupción generada. & \texttt{0} \\
    \hline
    \texttt{19--30[1]} & \texttt{P0.nFEI} & Igual que la fila anterior para los bits altos disponibles en Port 0. & \texttt{0} \\
    \hline
    \texttt{14:12} & \texttt{-} & Reservado. & NA \\
    \hline
    \texttt{31} & \texttt{-} & Reservado. & NA \\
    \hline
  \end{tabularx}
\end{table}

\begin{table}[H]
  \centering
  \small
  \caption{GPIO Interrupt Status for port 2 Falling Edge Interrupt \texttt{IO2IntStatF} (Table 121)}
  \begin{tabularx}{0.96\textwidth}{|>{\centering\arraybackslash}p{0.14\textwidth}|>{\raggedright\arraybackslash}p{0.18\textwidth}|>{\raggedright\arraybackslash}X|>{\centering\arraybackslash}p{0.10\textwidth}|}
    \hline
    \textbf{Bit} & \textbf{Symbol} & \textbf{Value / Description} & \textbf{Reset} \\
    \hline
    \texttt{0--11} & \texttt{P2.nFEI} & Estado de interrupción por flanco descendente en \texttt{P2.n}. \texttt{0}: no detectado. \texttt{1}: interrupción generada. & \texttt{0} \\
    \hline
    \texttt{12--13[1]} & \texttt{P2.nFEI} & Igual que la fila anterior para \texttt{P2.12} y \texttt{P2.13}. & \texttt{0} \\
    \hline
    \texttt{31:14} & \texttt{-} & Reservado. & NA \\
    \hline
  \end{tabularx}
\end{table}

\begin{table}[H]
  \centering
  \small
  \caption{GPIO Interrupt Clear register for port 0 \texttt{IO0IntClr} (Table 122)}
  \begin{tabularx}{0.96\textwidth}{|>{\centering\arraybackslash}p{0.14\textwidth}|>{\raggedright\arraybackslash}p{0.18\textwidth}|>{\raggedright\arraybackslash}X|>{\centering\arraybackslash}p{0.10\textwidth}|}
    \hline
    \textbf{Bit} & \textbf{Symbol} & \textbf{Value / Description} & \textbf{Reset} \\
    \hline
    \texttt{0--11, 15--18} & \texttt{P0.nCI} & Limpia interrupciones GPIO para \texttt{P0.n}. \texttt{0}: \texttt{IOxIntStatR/F} sin cambios. \texttt{1}: limpia los bits correspondientes. & \texttt{0} \\
    \hline
    \texttt{19--30[1]} & \texttt{P0.nCI} & Igual que la fila anterior para los bits altos disponibles en Port 0. & \texttt{0} \\
    \hline
    \texttt{14:12} & \texttt{-} & Reservado. & NA \\
    \hline
    \texttt{31} & \texttt{-} & Reservado. & NA \\
    \hline
  \end{tabularx}
\end{table}

\begin{table}[H]
  \centering
  \small
  \caption{GPIO Interrupt Clear register for port 2 \texttt{IO2IntClr} (Table 123)}
  \begin{tabularx}{0.96\textwidth}{|>{\centering\arraybackslash}p{0.14\textwidth}|>{\raggedright\arraybackslash}p{0.18\textwidth}|>{\raggedright\arraybackslash}X|>{\centering\arraybackslash}p{0.10\textwidth}|}
    \hline
    \textbf{Bit} & \textbf{Symbol} & \textbf{Value / Description} & \textbf{Reset} \\
    \hline
    \texttt{0--11} & \texttt{P2.nCI} & Limpia interrupciones GPIO para \texttt{P2.n}. \texttt{0}: \texttt{IOxIntStatR/F} sin cambios. \texttt{1}: limpia los bits correspondientes. & \texttt{0} \\
    \hline
    \texttt{12--13[1]} & \texttt{P2.nCI} & Igual que la fila anterior para \texttt{P2.12} y \texttt{P2.13}. & \texttt{0} \\
    \hline
    \texttt{31:14} & \texttt{-} & Reservado. & NA \\
    \hline
  \end{tabularx}
\end{table}

\begin{itemize}
  \item \textbf{[1]} No disponible en encapsulado de 80 pines.
\end{itemize}

\subsection{GPIO usage notes}
\subsubsection{Example: instantaneous output of 0s and 1s on a GPIO port}
\begin{itemize}
  \item El periférico permite actualización instantánea de patrones de salida mediante accesos de \textbf{32 bits}, \textbf{16 bits} o \textbf{8 bits}.
  \item La elección del ancho de acceso depende de la agrupación de señales y del patrón a escribir.
  \item El acceso de menor ancho simplifica cambios locales; el acceso de palabra completa simplifica patrones paralelos.
\end{itemize}

\subsubsection{Writing to \texttt{FIOSET}/\texttt{FIOCLR} vs. \texttt{FIOPIN}}
\begin{itemize}
  \item \texttt{FIOSET} y \texttt{FIOCLR} permiten cambiar bits puntuales sin reescribir el resto del puerto.
  \item Escribir en \texttt{FIOPIN} aplica un valor completo sobre todos los bits de salida no enmascarados.
  \item Para cambios puntuales de salida, la vía normal es \texttt{FIOSET}/\texttt{FIOCLR}.
  \item \texttt{FIOPIN} conviene cuando se desea forzar en una sola operación un patrón paralelo completo.
\end{itemize}
